\documentclass[12pt]{article}
\usepackage{fullpage}
\usepackage{amsmath,amssymb,amsthm}

\newtheorem{lemma}{Lemma}
\newtheorem{proposition}{Proposition}

\DeclareMathOperator{\Lk}{Lk}
\newcommand{\bbarotimes}{\mathbin{\overline{\otimes}}}
\newcommand{\cA}{\mathcal A}
\newcommand{\cH}{\mathcal H}
\newcommand{\cK}{\mathcal K}
\newcommand{\C}{\mathbb C}

\begin{document}

\title{Proper proximality of irreducible graph-product von Neumann algebras}
\author{}
\date{}
\maketitle

\section*{Problem statement and interpretation}

I use the standard graph-product meaning of irreducible: the finite graph
\(\Gamma=(V,E)\) is not a non-trivial join, equivalently the complement
\(\Gamma^c\) is connected.  The traces are faithful normal traces.  For
\(U\subset V\), \(M_U\) denotes the graph product over the induced subgraph and
\(E_U:M\to M_U\) its canonical expectation.

Put \(A_N=\mathbb B(L^2N)\) and
\[
        \cA_N=(A_N^{\sharp_J})^*
\]
for the simultaneous \(N\)- and \(JNJ\)-normal bidual of
Ding--Kunnawalkam Elayavalli--Peterson (DKEP) \cite{DKEP}.  If \(X\) is an
\(N\)-boundary piece, \(q_X\in\cA_N\) denotes the support projection of the
corner \((K_X(N)^{\sharp_J})^*\).  Thus
\[
 \widetilde S_X(N)=
 \{T\in\cA_N:[T,a]\in q_X\cA_Nq_X\hbox{ for all }a\in JNJ\}.
\]
For \(X=\mathbb K(L^2N)\) write \(q_0\) and \(\widetilde S(N)\).  We use DKEP
Lemma 8.5 for the trivial group action: if \(N\) is a separable finite factor,
then proper proximality relative to \(X\) is equivalent to the absence of an
\(N\)-central state on \(\widetilde S_X(N)\) restricting to the trace.

\section*{Normal supports and graph-product boundary corners}

Let \(A=\mathbb B(L^2N)\).  If \(D\subset A\) is a hereditary \(C^*\)-subalgebra
invariant under the multiplier actions of \(N\) and \(JNJ\), let
\(s_D\in A^{**}\) be its open support.  In each one-normal category, the
paragraph before DKEP Remark 2.8 identifies the normal bidual of \(D\) with
\(s_D(A^\sharp)^*s_D\).  DKEP Proposition 2.10 identifies simultaneous normal
biduals with intersections of the two one-normal biduals; hence
\[
        (D^{\sharp_J})^*=\cA_N\cap s_DA^{**}s_D                  \tag{1}
\]
inside \(A^{**}\).  This is a corner of \(\cA_N\); denote its unit by \(q(D)\).
Thus \(T\in q(D)\cA_Nq(D)\) exactly when \(T\in\cA_N\) and \(T=s_DTs_D\) in
\(A^{**}\).

For a DKEP boundary piece \(X\), \(q_X=q(K_X(N))\), but this is also \(q(X)\).
Indeed, if \(\phi\in A^{\sharp_J}\) vanishes on \(X\), then, since \(XAX\) is
norm dense in the hereditary algebra \(X\), DKEP Theorem 5.6 with \(X=Y\) gives
\(\phi(T)=0\) for every \(T\in K_X^{\infty,1}(N)\), hence for every
\(T\in K_X(N)\).  The converse is trivial from \(X\subset K_X(N)\).  Thus \(X\)
and \(K_X(N)\) have the same annihilator in \(A^{\sharp_J}\) and the same
weak-* closed corner in \(\cA_N\).

We shall use the following support consequence: if \(B,C,D\) are hereditary
subalgebras as above and
\[
        s_C(s_BA^{**}s_B)s_C\subset s_DA^{**}s_D,                \tag{2}
\]
then
\[
        q(C)\bigl(q(B)\cA_Nq(B)\bigr)q(C)\subset q(D)\cA_Nq(D).   \tag{3}
\]
This follows from (1).  We also use the elementary rectangular fact
\[
R^*R\in sA^{**}s,\quad S^*S\in tA^{**}t
       \quad\Longrightarrow\quad R^*S\in sA^{**}t,               \tag{4}
\]
because \(R=Rs\) and \(S=St\).

The next lemma supplies the coefficient-boundary identification used below.
For a von Neumann subalgebra \(B\subset N\), set
\[
        \cH_B={}_NL^2N_B\otimes_B{}_BL^2N_N .
\]

\begin{lemma}[induced subalgebra boundary]\label{lem:coeff}
For every expected subalgebra \(B\subset N\),
\[
        X_{\cH_B}=X_B ,
\]
where the left hand side is DKEP's coefficient boundary of the \(N\)-\(N\)
correspondence \(\cH_B\), and \(X_B\) is the subalgebra boundary of DKEP
Example 3.4.  In the tracial case \(\overline{\cH_B}\cong\cH_B\) as an
\(N\)-\(N\) correspondence.
\end{lemma}

\begin{proof}
Let \(H={}_NL^2N_B\) and \(K={}_BL^2N_N\).  DKEP Example 5.11 gives
\(X_H=X_B\).  On the \(B\)-side, the coefficient boundary is all
\(\mathbb B(L^2B)\), since for the vector \(\widehat 1\) the coefficient
\(T_{\widehat1}T_{\widehat1}^{\,*}\) (for \(H\)) and
\(T_{\widehat1}^{\,*}T_{\widehat1}\) (for \(K\)) is \(1_{L^2B}\).  Hence
\(H\) is mixing relative to \(X_B\times\mathbb B(L^2B)\), and \(K\) is mixing
relative to \(\mathbb B(L^2B)\times\mathbb B(L^2N)\).  DKEP Theorem 5.10 and
the fusion identity \(T_{\xi\otimes\eta}=T_\eta T_\xi\) show this explicitly:
if \(T_\xi\in\mathbb B(L^2N,L^2B)X_B\) and \(\eta\) is bounded in \(K\), then
\(T_{\xi\otimes\eta}\in\mathbb B(L^2N)X_B\), and such elementary tensors are
cyclic.  Thus \(H\otimes_BK=\cH_B\) is mixing relative to
\(X_B\times\mathbb B(L^2N)\).  By the minimality assertion in DKEP Example 5.11,
\(X_{\cH_B}\subset X_B\).

Conversely, if \(\xi,\eta\) are bounded vectors in \(H\), then
\(\xi\otimes\widehat1\) and \(\eta\otimes\widehat1\) are bounded in
\(\cH_B\), and
\[
 T_{\xi\otimes\widehat1}^{\,*}T_{\eta\otimes\widehat1}
 =T_\xi^{\,*}(T_{\widehat1}^{\,*}T_{\widehat1})T_\eta
 =T_\xi^{\,*}T_\eta .
\]
Thus every coefficient generator of \(X_H=X_B\) belongs to \(X_{\cH_B}\), so
\(X_B\subset X_{\cH_B}\).  Finally, the unitary
\(\overline{\xi\otimes\eta}\mapsto J\eta\otimes J\xi\) identifies
\(\overline{\cH_B}\) with \(\cH_B\).
\end{proof}

If \(Z\subset U\subset V\), then \(e_{M_Z}\le e_{M_U}\); Lemma
\ref{lem:coeff} and heredity give \(X_{M_Z}\subset X_{M_U}\).  We shall also
use the following immediate consequence.  If a correspondence is a direct sum of
correspondences \(\cK_j\) with \(X_{\cK_j}\subset X\), then all its coefficients
lie in \(s_XA^{**}s_X\): for finite sums this is (4) applied to the left ideal
of \(X\), and arbitrary direct sums follow by weak-* limits of finite
coordinate cutdowns.

\begin{lemma}[boundary-corner intersection]\label{lem:intersection}
Let \(q_U=q_{X_{M_U}}\).  For all \(U,W\subset V\),
\[
        q_W(q_U\cA_Mq_U)q_W\subset q_{U\cap W}\cA_Mq_{U\cap W}.   \tag{5}
\]
Consequently, if \(U_1,\ldots,U_m\subset V\) and \(I=\cap_iU_i\), then
\[
q_{U_m}q_{U_{m-1}}\cdots q_{U_1}q_{U_2}\cdots q_{U_{m-1}}q_{U_m}
\in q_I\cA_Mq_I .
\]
\end{lemma}

\begin{proof}
Let \(s_U\) be the ordinary support of \(X_{M_U}\) in \(A_M^{**}\), and put
\(I=U\cap W\).  By (3), it suffices to prove
\[
        s_W(s_UA_M^{**}s_U)s_W\subset s_IA_M^{**}s_I.             \tag{6}
\]

By Lemma \ref{lem:coeff}, \(X_{M_U}\) is the coefficient boundary of
\(\cH_U=L^2M\otimes_{M_U}L^2M\), and \(\overline{\cH_U}\cong\cH_U\).  Let
\(\mathcal C_U\) be the linear span of coefficients of all finite Connes fusions
of copies of \(\cH_U\) and \(\overline{\cH_U}\).  The identities
\(T_{\bar\xi}=T_\xi^*\) and \(T_{\xi\otimes\eta}=T_\eta T_\xi\) make
\(\mathcal C_U\) a self-adjoint algebra containing the positive coefficient
generators of \(X_{M_U}\); off-diagonal coefficients are in the hereditary
algebra generated by these positives by (4).  Therefore \(X_{M_U}\) has a
contractive approximate
unit whose elements are norm limits of elements of \(\mathcal C_U\) (take
\(f_\varepsilon(\sum_{j=1}^n T_{\xi_j}^*T_{\xi_j})\), \(f_\varepsilon(0)=0\),
and approximate by polynomials).

Choose such approximate units \(h_\alpha^W\) and \(h_\beta^U\).  We claim
\[
        h_\alpha^Wh_\beta^U\in s_IA_M^{**}s_U,\qquad
        h_\beta^Uh_\alpha^W\in s_UA_M^{**}s_I .                  \tag{7}
\]
It is enough, by norm closure, to check this for finite-fusion coefficients.
Let \(a=T_{\xi_1}^*T_{\xi_2}\in\mathcal C_W\) and
\(b=T_{\eta_1}^*T_{\eta_2}\in\mathcal C_U\).  Then
\[
        ab=T_{\xi_1\otimes\overline{\xi_2}\otimes\eta_1}^{\,*}T_{\eta_2}.
                                                                        \tag{8}
\]
CDHJKN identify \(\cH_U\) with \(H_U(V,V)\), and their fusion rule
\cite[Theorem 5.4 and Proposition 5.5]{CDHJKN}, applied repeatedly, decomposes
the correspondence containing
\(\xi_1\otimes\overline{\xi_2}\otimes\eta_1\) as a direct sum of copies of
\(\cH_Z\) with \(Z\subset I\).  The direct-sum support observation and
\(X_{M_Z}\subset X_{M_I}\) imply that the first vector in (8) has diagonal
coefficient supported by \(s_I\); the finite fusion containing \(\eta_2\) similarly decomposes over subsets of
\(U\), so \(\eta_2\) has diagonal coefficient supported by \(s_U\).  The
rectangular support fact (4) gives the first
inclusion in (7), and the second is the adjoint.

For \(x\in s_UA_M^{**}s_U\), the weak-* limits
\(h_\alpha^W\to s_W\) and \(h_\beta^U\to s_U\) give
\[
h_\alpha^Wh_\beta^U\,x\,h_\gamma^Uh_\delta^W\in s_IA_M^{**}s_I,
\]
and successive weak-* limits yield (6).  Thus (5) follows from (3).  The
palindromic assertion follows by induction from (5).
\end{proof}

\section*{A finite boundary criterion}

\begin{lemma}\label{lem:finite}
Let \(N\) be a separable non-amenable finite factor.  Let \(X_1,\ldots,X_m\) be
boundary pieces and put \(q_i=q_{X_i}\).  Assume
\[
        q_mq_{m-1}\cdots q_1q_2\cdots q_{m-1}q_m\in q_0\cA_Nq_0.  \tag{9}
\]
If \(N\) is properly proximal relative to each \(X_i\), then \(N\) is properly
proximal.
\end{lemma}

\begin{proof}
Assume \(N\) is not properly proximal.  DKEP Lemma 8.5 gives an \(N\)-central
state \(\omega\) on \(\widetilde S(N)\) with \(\omega|_N=\tau\).  Fix \(i\) and
write \(q=q_i\).  Since \(q\) commutes with \(N\) and \(JNJ\),
\(q^\perp\widetilde S_{X_i}(N)q^\perp\subset\widetilde S(N)\).  If
\(\omega(q^\perp)>0\), compression by \(q^\perp\) gives an \(N\)-central state
on \(\widetilde S_{X_i}(N)\).  Its restriction to \(N\) is a normal central
positive functional, hence a scalar multiple of \(\tau\) because \(N\) is a
factor; after normalization it restricts to \(\tau\), contradicting relative
proper proximality.  Thus \(\omega(q_i^\perp)=0\) for every \(i\).

Extend \(\omega\) to a state on \(\cA_N\).  In the GNS representation every
\(q_i\) fixes the cyclic vector, so the positive contraction in (9) has
\(\omega\)-value \(1\).  Since it lies in \(q_0\cA_Nq_0\), we get
\(\omega(q_0)=1\).  DKEP's canonical \(N\)- and \(JNJ\)-bimodular u.c.p. map
\(\Theta:\mathbb B(L^2N)\to q_0\cA_Nq_0\) sends \(1\) to \(q_0\) and is the
identity on \(N\) after compression.  Hence \(\omega\circ\Theta\) is a Connes
hypertrace on \(\mathbb B(L^2N)\), contradicting non-amenability of \(N\).
\end{proof}

\section*{A relative Bass--Serre paradox}

\begin{proposition}\label{prop:BS}
Let \(N=P*_{B}Q\) with trace-preserving expectations and suppose
\(Q=B\bbarotimes A\).  Assume that \(A\) contains a trace-zero unitary \(a\) and
that \(P\) contains unitaries \(w_1,w_2\) satisfying
\[
        E_B(w_i)=0,\qquad E_B(w_1^*w_2)=0 .
\]
Then \(N\) is properly proximal relative to \(X_B\).
\end{proposition}

\begin{proof}
Let \(H_P=L^2P\ominus L^2B\) and \(H_Q=L^2Q\ominus L^2B\).  In the AFP
decomposition of \(L^2N\), let \(p\) and \(q\) be the projections onto reduced
words whose first letter lies in \(H_P\) and \(H_Q\), respectively; then
\(1=e_B+p+q\).  For \(R_x=Jx^*J\), reduced-word inspection gives
\[
[p,R_x]=R_xe_B-e_BR_x,\ [q,R_x]=0\quad(x\in P\ominus B),
\]
and the same formula with \(p,q\) interchanged for \(x\in Q\ominus B\); elements
of \(B\) commute with \(p,q\).  These commutators are \(N\)- and \(JNJ\)-
translates of \(e_B\), so DKEP Lemma 6.1 implies \(p,q\in S_{X_B}(N)\).

Let \(\varphi\) be an \(N\)-central state on \(S_{X_B}(N)\) normal on \(N\).
Left multiplication by a centered \(P\)-letter sends \(qL^2N\) into \(pL^2N\),
and \(q w_1^*w_2q=0\) because \(E_B(w_1^*w_2)=0\).  Hence the two ranges are
orthogonal and
\(w_1qw_1^*+w_2qw_2^*\le p\).  Similarly, since \(a\in Q\ominus B\),
\(apa^*\le q\).  Centrality gives
\[
        2\varphi(q)\le\varphi(p)\le\varphi(q),
\]
so \(\varphi(p)=\varphi(q)=0\).  The unitary \(h=w_1a\) is \(B\)-Haar: every
non-zero power is a reduced alternating word and has \(B\)-expectation zero.
Thus the projections \(h^ne_Bh^{-n}\) are pairwise orthogonal and have the same
\(\varphi\)-value; hence \(\varphi(e_B)=0\).  This contradicts
\(1=e_B+p+q\).  DKEP Theorem 6.2 gives proper proximality relative to \(X_B\).
\end{proof}

\section*{Application to graph products}

Fix \(v\in V\) and let \(L_v=\Lk_\Gamma(v)\).  The graph product splits as
\[
        M\cong M_{V\setminus\{v\}}*_{M_{L_v}}
        (M_{L_v}\bbarotimes M_v).                                \tag{10}
\]
Let \(C(v)=\{r\ne v:r\not\sim_\Gamma v\}\).  Since \(\Gamma^c\) is connected,
\(C(v)\ne\emptyset\).  If \(C(v)\) contains distinct \(r,t\), choose trace-zero
unitaries \(u_r\in M_r,u_t\in M_t\) and put \(w_1=u_r,w_2=u_t\).  If
\(C(v)=\{r\}\), connectedness and \(|V|\ge3\) give
\(s\in L_v\) with \(s\not\sim_\Gamma r\); put \(w_1=u_r,w_2=u_su_r\).  In both
cases reduced-word calculus gives
\[
        E_{M_{L_v}}(w_i)=0,\qquad E_{M_{L_v}}(w_1^*w_2)=0 .
\]
The trace-zero unitary in \(M_v\) is the unitary \(a\) in (10), so
Proposition \ref{prop:BS} shows that \(M\) is properly proximal relative to
\(X_{M_{L_v}}\) for every \(v\in V\).

Enumerate \(V=\{v_1,\ldots,v_m\}\).  Since no vertex belongs to its own link,
\(\cap_iL_{v_i}=\emptyset\).  Lemma \ref{lem:intersection} gives
\[
 q_{L_{v_m}}\cdots q_{L_{v_1}}q_{L_{v_2}}\cdots q_{L_{v_m}}
 \in q_{\emptyset}\cA_Mq_{\emptyset}.
\]
Here \(M_\emptyset=\C\), and \(X_{M_\emptyset}=\mathbb K(L^2M)\).

Finally, CDHJKN Main Theorem 0.5 implies that \(M\) is a factor: the possible
obstructions would give either a universal vertex of \(\Gamma\) or an isolated
exceptional non-edge component, both impossible because \(\Gamma^c\) is
connected with at least three vertices.  CDHJKN Proposition 6.3 similarly
excludes amenability: its amenable non-edge components are single vertices or
isolated exceptional edges, whereas \(\Gamma^c\) is connected with at least
three vertices.  Thus \(M\) is a separable non-amenable finite factor.
Lemma \ref{lem:finite} applies and proves that \(M\) is properly proximal.

\begin{thebibliography}{9}

\bibitem{DKEP}
C.~Ding, S.~Kunnawalkam Elayavalli, and J.~Peterson,
\emph{Properly proximal von Neumann algebras},
Duke Math. J. \textbf{172} (2023), 2821--2894.

\bibitem{CDHJKN}
I.~Charlesworth, R.~de Santiago, B.~Hayes, D.~Jekel,
S.~Kunnawalkam Elayavalli, and B.~Nelson,
\emph{On the structure of graph product von Neumann algebras},
Publ. Res. Inst. Math. Sci. \textbf{61} (2025), 713--762.

\end{thebibliography}

\end{document}
